\documentclass[11pt,a4paper]{extarticle}

\usepackage[T2A]{fontenc}
\usepackage[utf8]{inputenc}
\usepackage[russian]{babel}
\usepackage{indentfirst}
\usepackage{geometry}
\usepackage{setspace}

\geometry{left=20mm, right=10mm, top=15mm, bottom=15mm}
\onehalfspacing
\setlength{\parindent}{1.25cm}
\setlength{\parskip}{0pt}
\frenchspacing

\begin{document}
\pagestyle{empty}

\begin{center}
  \textbf{РЕЦЕНЗИЯ}\\
  \textbf{на выпускную квалификационную работу}\\
  Киселева Никиты Сергеевича\\
  <<Методы определения достаточного размера выборки\\ в параметрических моделях>>
\end{center}

\vspace{0.3em}

Выпускная квалификационная работа Н.\,С.~Киселева посвящена количественному анализу стабилизации поверхности эмпирической функции потерь параметрических моделей при увеличении объема обучающей выборки и построению критериев достаточного размера выборки. Тема работы актуальна: существующие критерии стабилизации ландшафта (одноточечный и изотропный среднеквадратичный) практически неприменимы в моделях современных размеров из-за зависимости доминирующей константы их верхних оценок от размерности пространства параметров как~${N^2}$. Постановка задачи поиска критериев, согласованных с эмпирически наблюдаемой анизотропией спектра матрицы Гессе, своевременна и хорошо мотивирована. Достоверность результатов обеспечивается строгостью доказательств в рамках четко обозначенных предположений, воспроизводимостью экспериментов и согласием полученных численных показателей с предсказанными теоретическими оценками; программная реализация доступна публично.

Теоретическая значимость работы состоит в построении унифицированной схемы критериев стабилизации с произвольной функцией предпочтения точек в пространстве параметров и в доказательстве для подпространственного критерия~${\Delta_2^{(D)}}$ верхней оценки~${\mathcal O(k^{-2})}$, в которой зависимость от полной размерности заменена на зависимость от размерности подпространства~${D}$. Практическая значимость заключается в количественном инструменте для выбора достаточного объема обучающей выборки и в масштабируемых алгоритмах, обеспечивающих ускорение примерно в~${10^4}$ раз по сравнению с прямой Монте-Карло оценкой без потери точности в локальном режиме применимости.

К положительным сторонам работы следует отнести единство теоретического анализа и эмпирической проверки на двух моделях существенно разного масштаба (полносвязной сети на MNIST и трансформере объемом~${\sim 10^8}$ параметров), академическую аккуратность формулировок с явным указанием идеализирующих предположений, логически связную структуру изложения и апробацию результатов на рецензируемых научных площадках. В качестве замечания можно отметить, что эмпирическая проверка работоспособности подпространственного критерия в моделях большой размерности проведена на одной трансформерной архитектуре (\texttt{nanochat~d6}, ${\sim 10^8}$ параметров) и одной контрольной точке обучения. Расширение экспериментов на другие семейства декодерных моделей и на различные стадии обучения дополнительно укрепило бы эмпирическую часть работы. Указанное ограничение, однако, оправдано вычислительной стоимостью повторных операций второго порядка в моделях такого размера и не снижает значимости теоретического вклада, носящего методологический характер.

Работа в полной мере соответствует требованиям, предъявляемым к выпускным квалификационным работам магистров в МФТИ, заслуживает оценки \textbf{<<отлично>>}, а её автор рекомендуется к продолжению научной деятельности в аспирантуре.

\vspace{0.5em}

\begin{flushleft}
\hfill\begin{tabular}{l}
\textbf{Рецензент}:\\
Ивахненко Андрей Александрович, к.ф.-м.н.,\\
младший научный сотрудник Лаборатории №42\\
Института проблем управления РАН им. Трапезникова
\end{tabular}
\end{flushleft}

\vfill

\begin{flushright}
    <<\underline{\hspace{1cm}}>>\ \underline{\hspace{2.5cm}}\ 2026~г.
\end{flushright}

\end{document}
